% Part 0: the complete statement, the roadmap and the logical skeleton.
% ASCII source only. \input after page 19 of report.tex.
\proofpart{0}
\pagetitle{20 / THE COMPLETE STATEMENT}{One conjecture, held up by eight layers}
{\small\color{teal}\textbf{Added 27 November 2026.} The front of the theorem chain. Parts I--V follow.}\par\smallskip

\textbf{What follows, and what it is not.} Everything in Parts I--V is one
construction. It is organised by the structures that every principal bundle
carries, taken as layers: the \emph{base} (space, its size, and the scale
along which it is refined), the \emph{structure group} $SU(2)$ with its centre,
the \emph{connection} and its holonomies, the \emph{gauge quotient}
$\mathcal A/\mathcal G$, its \emph{topology}, the \emph{measure} on the quotient,
the \emph{spectrum} of the Hamiltonian, and the \emph{scale fibration} that
splits slow modes from fast ones. Each result lives in one layer, and each
implication between layers is marked with what holds it up. The chain is
complete in the sense that every step the Main Conjecture needs is named, so
the whole can be checked and can be wrong. It is
not validated, and the Main Conjecture is not proved. The proved pieces are standard mathematics, or
theorems about a lattice, two dimensions, or a finite-dimensional reduction;
none is a statement about four-dimensional continuum Yang--Mills.

\textbf{Two words of jargon.} \emph{Ultraviolet} means the fine-grain
direction: zooming in, where the field has wiggles at ever shorter distances,
infinitely many of them, and the lattice spacing goes to zero. \emph{Infrared}
means the opposite direction: zooming out, as the room grows without end. (The
names are borrowed from light only by analogy: short wavelengths are violet,
long ones red.) In this programme the ultraviolet is the \emph{fast modes} of
the scale fibration (Part IV, M3) and the infrared is \emph{the room growing}
(Part V, M2).

\textbf{Conventions.} $G=SU(2)$. The coupling runs at one loop as
$1/g^2(\mu)=(b_0/8\pi^2)\ln(\mu/\Lambda)$ with $b_0=22/3$, for the action
$\frac{1}{4g^2}\int F^a_{\mu\nu}F^{a\,\mu\nu}$. The scale $\Lambda$ is generated
by the running and does not appear in the classical theory. Its value depends on
the scheme; the ratio between two schemes is a fixed number, so ``$\Delta=c\Lambda$''
means that $c$ is a pure number once the scheme is fixed.

\begin{mainconj}
\status{CONJECTURE}
\begin{enumerate}[label=(\roman*),leftmargin=18pt,itemsep=2pt,topsep=2pt]
\item \emph{(Existence and the gap.)} There is a probability measure on
(generalised) $SU(2)$ connections on $\mathbb R^4$ modulo gauge, obtained as a
limit of Wilson lattice measures as the spacing $a\to0$ along the running
coupling and the volume $L^4\to\infty$, whose Schwinger functions of
gauge-invariant observables (smeared $\operatorname{tr}F^2$ and smeared Wilson
loops) satisfy the Osterwalder--Schrader axioms: Euclidean invariance,
reflection positivity, regularity, and clustering. By the
Osterwalder--Schrader reconstruction theorem this gives a Wightman theory with
Hilbert space $\mathcal H$, vacuum $\Omega$ and Hamiltonian $H\geq0$. The theory
is non-trivial, which I take to mean that its truncated four-point function
does not vanish identically, so it is not a generalised free field. It has a mass gap:
$\operatorname{spec}(H)\subset\{0\}\cup[\Delta,\infty)$ with $0$ a simple
eigenvalue, and $\Delta=c\,\Lambda$ with $c>0$. In particular the truncated
correlations of separated gauge-invariant probes decay like
$e^{-\Delta\cdot\mathrm{distance}}$.
\item \emph{(The sharper form: the note is inherited.)} The lowest note has a
traceable carrier along the whole direction of scale.
\begin{itemize}[leftmargin=14pt,itemsep=1pt,topsep=1pt]
\item At the lattice end (Hamiltonian lattice theory, strong coupling) the
lowest gauge-invariant excitation is, at leading order, one plaquette of
$j=\tfrac12$ flux, with energy $4\cdot\frac{g^2}{2a}\cdot C_2(\tfrac12)=\frac{3g^2}{2a}$,
where $C_2(\tfrac12)=\tfrac34$ is the Casimir of one turn.
\item In the small room (homogeneous connections on the Berger sphere, window
$t>1$, $1<s<(t+2)/(t+1)$) the lowest note is carried by Helike, the
charge-$2$ maximal-charge twist, for $s>s^*(g^2)$, and by the charge-$0$
competitor for $s<s^*(g^2)$, with a single identifiable switch at $s^*$.
\item Between and beyond these, each carrier keeps a measurable property (in the
reduction, an own-block share $\geq\tfrac12$; on the lattice, representation
content, loop extent and Wilson-loop overlap), which changes hands only at
identifiable level crossings, uniformly as the cutoffs are removed; and the gap
in the field's own scale, $\Delta/\Lambda$, stays bounded away from zero.
\end{itemize}
\end{enumerate}
\end{mainconj}

Part (i) is the Jaffe--Witten statement of the problem (Clay Mathematics
Institute, 2000), restricted to $SU(2)$; the axioms are those of Osterwalder
and Schrader (Commun.\ Math.\ Phys.\ 31 (1973) 83; 42 (1975) 281). Part (ii) is
mine. It is stronger than (i), and it can fail while (i) holds: a gap made at
weak coupling with no inherited carrier would refute (ii) and leave (i)
standing.

\begin{sketch}
The argument is the logical skeleton below, read top to bottom. The key
estimate is a bound, uniform in the ultraviolet cutoff and in the size of the
room, on how far the fast modes and the larger rooms can move the lowest
levels of the reduction relative to the gap between the carriers (Parts IV and
V). It can fail in three places, each named in advance: the stated hypothesis
$S_C$ (Part III), a carrier that dissolves without a switch (the M8 and M9
tests), and $\Delta R$ flattening as the room grows (M2). Even if all three
hold, the passage from bounds on rooms to Osterwalder--Schrader limits is the
constructive problem itself, and nothing here does it.
\end{sketch}

\textbf{The five tags.} Every result in Parts I--V carries exactly one.
\begin{center}\small
\begin{tabular}{@{}p{1.55in}p{4.75in}@{}}\toprule
\textbf{Tag} & \textbf{Meaning}\\\midrule
\status{PROVED HERE} & A complete proof is written in this report.\\[2pt]
\status{PROVED, CITED} & Proved in the cited work; not reproved here.\\[2pt]
\status{COMPUTED} & A numerical result, stated with its number and with what it checks. It is evidence, not a proof.\\[2pt]
\status{CONDITIONAL ON ...} & Proved, assuming the named hypothesis, which is itself open.\\[2pt]
\status{CONJECTURE} & Not proved. Followed by my best-guess argument and where it could fail.\\\bottomrule
\end{tabular}
\end{center}

\textbf{Roadmap.} Counts refer to the principal results named in this table
(PH proved here, PC proved and cited, C computed, Cd conditional, X
conjecture), counted from the parts as written. Each part's own tags are
authoritative.
\begin{center}\small
\begin{tabular}{@{}p{0.55in}p{1.25in}p{3.45in}p{0.95in}@{}}\toprule
\textbf{Part} & \textbf{Layers} & \textbf{Main results} & \textbf{Status}\\\midrule
I. The bundle and the lattice & base (lattice), structure group, connection, measure & holonomy is gauge covariant and Wilson loops are gauge invariant, in the continuum and on the lattice; loops determine the connection (Giles); the Wilson action (Wilson); the lattice measure lives on the quotient; why the quotient (Gribov, Singer); $SU(2)$ as the double cover, its centre and Casimirs; reflection positivity and the transfer matrix, and the strong-coupling gap (Osterwalder--Seiler); the electric gap is exactly $3g^2/(2a)$; my reading of it as the leading term of the full gap & 6 PH, 5 PC, 1 X\\[4pt]
II. Two dimensions, exactly & connection, structure group & central averaging; the area law $\langle\chi_j\rangle/d_j=(I_{2j+1}/I_1)^{|A|}$ for every representation; the string tension is the Casimir, with Casimir scaling to $8/3$; centre vortices flip exactly the half-integer loops; the continuum two-dimensional theory (Migdal; Gross and Witten); the GIGI verifications & 5 PH, 1 PC, 1 C\\[4pt]
III. The small room & base (Berger spheres), gauge quotient, topology, spectrum & $\mathbb R^9/SO(3)$ with the orbit-volume measure; the Maurer--Cartan line, $CS(\lambda)=3\lambda^2-2\lambda^3$, the vacua, the background exactly $s\,q_{\mathrm{MC}}$; Helike (\emph{By Enough}; Floor v3.3) and her doublet in the reduced model; Theorem A; $S_C$; the ordering; the fold & 7 PH, 1 PC, 2 C, 1 Cd, 2 X\\[4pt]
IV. The scale fibration (M3) & scale fibration, base (running coupling) & the Berry connection is flat; the empty-sphere fibre; gauge equivalence of the flat points, so the regulator must be spectral; the divergence of the dressing is the one-loop running of $1/g^2$; the finite dressing, and at the background it widens Helike's lead; first-order shifts; the uniform-dressing conjecture & 6 PH, 2 C, 1 X\\[4pt]
V. Carriers, the room growing, the assembly & spectrum, topology, measure, base (scale) & the switch is an exact crossing; carrier persistence (M8) and the glueball link; $\theta$ as a Bloch phase and the barrier on the line; M2; the source bridge; partial ultraviolet results; reflection positivity survives limits; ultraviolet existence; the assembly & 4 PH, 1 PC, 1 Cd, 6 X\\[4pt]
VI. Where this stands & --- & each part against the established line it continues; where a referee will push & ---\\\bottomrule
\end{tabular}
\end{center}

\textbf{The logical skeleton.} Each arrow carries its status. The proofs are in
the parts named; they are not restated here.
\begin{enumerate}[leftmargin=18pt,itemsep=3pt,topsep=2pt]
\item \textbf{Part I $\Rightarrow$ a lattice theory.} On every finite lattice the
Wilson measure is a gauge-invariant probability measure that descends to the
finite quotient $\mathcal A/\mathcal G$, with no gauge fixing
\status{PROVED HERE}. It is reflection positive, so it has a transfer matrix and
a positive Hamiltonian, and for small $\beta$ it has a mass gap
(Osterwalder and Seiler, Ann.\ Phys.\ 110 (1978) 440) \status{PROVED, CITED}.
\item \textbf{Part I $\Rightarrow$ the note at the lattice end.} The electric
(strong-coupling) Kogut--Susskind Hamiltonian has a unique ground state and
first excited level exactly $4\cdot\frac{g^2}{2a}\cdot\frac34=\frac{3g^2}{2a}$,
carried by one plaquette in $j=\tfrac12$ \status{PROVED HERE}. That this is the
leading strong-coupling term of the gap of the full Hamiltonian
\status{CONJECTURE}.
\item \textbf{Part II $\Rightarrow$ the note, exactly, in two dimensions.} The
area law holds exactly for every representation, with the Casimir as the string
tension, and the centre vortices flip exactly the double-covered loops
\status{PROVED HERE}; GIGI reproduces them from point queries
\status{COMPUTED}.
\item \textbf{Part III $\Rightarrow$ a finite quotient with its topology.}
Homogeneous connections on the Berger sphere form $\mathbb R^9/SO(3)$; the
background sits on the Maurer--Cartan line at $\lambda=s$, just past the
winding-$1$ vacuum ($1<s<(t+2)/(t+1)$), with $CS=3s^2-2s^3$
\status{PROVED HERE}.
\item \textbf{Part III $\Rightarrow$ the carrier in the small room.} Helike is
an exact eigenvector of the Hessian on Berger spheres, with a closed-form
eigenvalue, and is its floor on the domain stated there (\emph{By Enough},
doi:10.5281/zenodo.24472820; Floor v3.3) \status{PROVED, CITED}; the
reduced Hamiltonian has discrete spectrum (Theorem A) \status{PROVED HERE};
Helike's level lies below the competitor's at three of four weak-coupling
controls \status{CONDITIONAL ON $S_C$}. $S_C$ itself: $C_3$, the fold, the
outside bound and the enclosed base are proved, the homotopy margin
$\ell_5-\rho_C=+0.21$ to $+0.24$ is floating point only \status{CONJECTURE}.
\item \textbf{Part V $\Rightarrow$ the switch.} The lowest note passes from
Helike to the competitor at one identifiable $s^*(g^2)$, with registered
predictions $1.1674$, $1.1809$, $1.1931$ at $g^2=0.03$, $0.05$, $0.07$; its
rise with $g^2$ is the small-volume precursor of $0^{++}<2^{++}$
\status{CONJECTURE}.
\item \textbf{Part IV $\Rightarrow$ the fast modes only dress.} The
$q$-dependent ultraviolet divergence of the fast zero-point energy is
$(b_0/8\pi^2)\operatorname{Vol}(S^3_t)$ times the classical energy density,
ratio $1.0000\pm10^{-4}$ at every $\lambda\in\{\tfrac18,\dots,\tfrac78\}$; at
the background the finite first-order dressing widens Helike's lead over the
competitor ($0.0165\to0.0191$ at $g^2=0.002$) \status{COMPUTED}. That the finite remainder is bounded uniformly in the cutoff
and does not exchange the carriers except at the switch (M3)
\status{CONJECTURE}.
\item \textbf{Part V $\Rightarrow$ the room grows.} The sources are removed or
matched to the unsourced theory, and $\Delta R/(R\Lambda)$ stays bounded away
from zero as $R\to\infty$ (M2); on the lattice the strong-coupling carrier
descends to weak coupling with only identifiable switches (M9)
\status{CONJECTURE}.
\item \textbf{Assembly (Part V).} Reflection positivity, a closed condition,
passes to any limit of the lattice expectations of bounded local observables,
for reflections in hyperplanes that are lattice hyperplanes along the sequence
\status{PROVED HERE}. Given steps
1--8, the limit $a\to0$, $L\to\infty$ exists, restores Euclidean invariance,
is non-trivial and has $\Delta=c\Lambda>0$, with the carriers of step 2 and
steps 4--5 as the inherited note: this is the Main Conjecture
\status{CONJECTURE}.
\end{enumerate}

\textbf{Where the weight rests.} Steps 1--4 are proved, cited or computed,
except my reading of $3g^2/(2a)$ as the leading term of the full lattice gap.
Step 5 rests on one hypothesis, $S_C$, whose remaining piece is an interval
Lehmann computation (R2) that is still running; until it closes, the ordering
is conditional. Steps 6--9 are conjecture apart from the computed parts of step 7, and step 9 is the whole constructive
problem: even granted steps 1--8, I cannot prove that the limit exists.
